\documentclass[12pt,superscriptaddress]{revtex4}
%\documentclass[prl,superscriptaddress,twocolumn]{revtex4}
%\usepackage[lite,subscriptcorrection,slantedGreek,nofontinfo]{mtpro2}
%\usepackage{times}
%\usepackage{CJK}
%\usepackage{ctex}

\usepackage{amsfonts}
\usepackage{mathrsfs}
\usepackage{amsmath}
\usepackage{amssymb}
%\usepackarrge{mathptmx}
\usepackage{bm}
\usepackage{braket}
\usepackage{feynmp-auto}
\usepackage[normalem]{ulem}
\usepackage{extarrows}
\usepackage{slashed}
\usepackage{isodateo}
\usepackage{graphicx}
\usepackage{xcolor}
\usepackage[CJKbookmarks=true, bookmarksnumbered=true,bookmarksopen=true]{hyperref}
\hypersetup{colorlinks,%
	linkcolor=[rgb]{0,0.2,0.6}, %
	citecolor=[rgb]{0,0.2,0.6},
	urlcolor=[rgb]{0,0.2,0.6}}

\usepackage{indentfirst}

\graphicspath{{fig/}}


\linespread{1.1}
\renewcommand{\FR}[2]{\displaystyle\frac{\,{#1}\,}{#2}}
\newcommand{\fr}[2]{\mbox{$\frac{\,{#1}\,}{#2}$}}

\renewcommand{\thefootnote}{\arabic{footnote}}

\newcommand{\order}[1]{\mathcal{O}({#1})}



\newcommand{\del}[1]{{\blue\sout{#1}}}


\newcommand{\red}{\color{red}}
\newcommand{\blue}{\color{blue}}

\newcommand{\nn}{\nonumber}
\newcommand{\nd}{\mathrm{d}}
\newcommand{\nD}{\mathrm{D}}
\newcommand{\nT}{\mathrm{T}}
\newcommand{\nW}{\mathrm{W}}
\newcommand{\nds}{\hat{\mathrm{d}}}
\newcommand{\mN}{\mathrm{N}}
\newcommand{\mK}{\mathcal{K}}
\newcommand{\mG}{\mathcal{G}}
\newcommand{\spdot}{\!\cdot\!}
\newcommand{\nG}{\mathrm{G}}
\newcommand{\ep}{\epsilon}
\newcommand{\w}{\omega}


\newcommand{\ob}[1]{\mkern 2mu \overline{\mkern -2mu #1 \mkern -2mu}\mkern 2mu}
\newcommand{\wt}[1]{\mkern 2mu \widetilde{\mkern -2mu #1 \mkern -2mu}\mkern 2mu}
\newcommand{\wh}[1]{\mkern 2mu \widehat{\mkern-2mu#1\mkern-2mu}\mkern 2mu}
\def\ma{\mathcal}





\begin{document}
	
	
	\title{Response to the Review Comments}
    \author{Jiaqi Chen}
    \author{Bo Feng}
	
	
	\maketitle

\noindent Dear Editors and Reviewers:
	
We would like to thank you for your careful and patient reading, helpful comments, and constructive suggestions, which have significantly improved the presentation of our manuscript. We have carefully considered all comments from the reviewers and revised our manuscript accordingly. The manuscript has also been double-checked, and the typos and grammar errors we found have been corrected. In the following section, we summarize our responses to these questions which need further explanation from the reviewers. We hope that our responses have well addressed all concerns from the reviewers.

\noindent\rule[0.25\baselineskip]{\textwidth}{1pt}


%--------------------------------------------
%---------------------------------------------

\begin{center}
\fcolorbox{black}{gray!10}
{
	\parbox{1\linewidth}
	{General Comments: on the confusion of the use of the abbreviation \textbf{CDE}.}
}
\end{center}
\noindent\textcolor{red}{Response}:
We reviewed all the ``CDE" in the paper and modified a part of them to ``CDE method" when referring to the method rather than the differential equations.


%--------------------------------------------
%---------------------------------------------

\begin{center}
\fcolorbox{black}{gray!10}
{
\parbox{1\linewidth}
{Page 1, Line 5: What is the function family in this context? Shouldn’t it be the integral family?}
}
\end{center}
\noindent\textcolor{red}{Response}:
We modified the ``function family" here to ``integral family."


%--------------------------------------------
%---------------------------------------------


\begin{center}
\fcolorbox{black}{gray!10}
{
	\parbox{1\linewidth}
	{		
Page 2, Line 21: Who is people? Citation needed. Also: It should
be master integrands not master integrals, when you speak about the forms.
	}
}
\end{center}
\noindent\textcolor{red}{Response}:
Originally, we cited the relevant articles at the end of this long sentence. To make it easier for the reader to notice, we moved the citation to an earlier position, near the ``people".


%--------------------------------------------
%---------------------------------------------

\begin{center}
\fcolorbox{black}{gray!10}
{
	\parbox{1\linewidth}
	{		
Page 2, Line 24: However, the origin of CDE is not clear. The origin of the method should be clear and cited...
	}
}
\end{center}
\noindent\textcolor{red}{Response}:
We modified it to ``However,  why and how such a structure could emerge from $\nd \log$-form integrand of master integral, is not fully understood.", to make it clear.

%--------------------------------------------
%---------------------------------------------

\begin{center}
\fcolorbox{black}{gray!10}
{
\parbox{1\linewidth}
  {		
  Page 4, Line 10: equivalence classes $\bra{\varphi}$ of \textbf{integrals} $\to$ ... \textbf{integrands} ...
  }
}
\end{center}
\noindent\textcolor{red}{Response}:
We modified it to ``integrands".

%--------------------------------------------
%---------------------------------------------


\begin{center}
\fcolorbox{black}{gray!10}
{
\parbox{1\linewidth}
  {		
  Page 4, Line 14: With [the] intersection number as the IBP invariant inner product $\to$ The intersection pairing is the product ...
  }
}
\end{center}
\noindent\textcolor{red}{Response}:
we modified the ``intersection number" to "the algorithm of intersection number", to make it clear that we mean the computation of intersection pairing rather than the number result.

%--------------------------------------------
%---------------------------------------------
\begin{center}
\fcolorbox{black}{gray!10}
{
\parbox{1\linewidth}
  {		
  Page 5: What happens between line 2 and 3 of (2.12)? Please elaborate (shortly).
  }
}
\end{center}
\noindent\textcolor{red}{Response}:
We added an explanation under equation (2.12), indicating that it is achieved through IBP.



%--------------------------------------------
%---------------------------------------------
\begin{center}
\fcolorbox{black}{gray!10}
{
\parbox{1\linewidth}
  {		
  Page 6, (2.18): The notation $\nabla^{-1} \varphi_L$ is not properly defined ...
  }
}
\end{center}
\noindent\textcolor{red}{Response}:
We added the explanation ``where any $g=\nabla_i^{-1}  f$ means the $g$ satisfies $\nabla_i g = f$." under equation (2.18). 


%--------------------------------------------
%---------------------------------------------
\begin{center}
\fcolorbox{black}{gray!10}
{
\parbox{1\linewidth}
  {		
  Page 12, Line 30: But we check our computation ... I do not understand what your check consisted of from this sentence. Please clarify.
  }
}
\end{center}
\noindent\textcolor{red}{Response}:
To make it clear, we modified it to ``To reconfirm our calculations of the intersection numbers (such as the computations in section 5) using the region method, we compute them again again by the fibration algorithm of multivariate intersection number. It is another algorithm that could compute multivariate intersection number, but does not suffer from the complication of the multivariate complex function."


%--------------------------------------------
%---------------------------------------------
\begin{center}
\fcolorbox{black}{gray!10}
{
\parbox{1\linewidth}
  {		
  Page 19, Line 19: ... in (2.87) should be [non?] zero
  }
}
\end{center}
\noindent\textcolor{red}{Response}:
Yes, thanks for indicate the typo.

%--------------------------------------------
%---------------------------------------------
\begin{center}
\fcolorbox{black}{gray!10}
{
\parbox{1\linewidth}
  {		
  Page 19, Line 27: By what are they ordered?
  }
}
\end{center}
\noindent\textcolor{red}{Response}:
We added the explanation after the sentence: ``(notice that $\varphi_a\sim \varphi_b, \varphi_b \sim \varphi_c$ do not give $\varphi_a \sim \varphi_c$, that is why there is a order)".


%--------------------------------------------
%---------------------------------------------
\begin{center}
\fcolorbox{black}{gray!10}
{
\parbox{1\linewidth}
  {		
  Page 19, last line: ``For example, for the case with seven master integrals" Which case? Refer to other place in the paper?
  }
}
\end{center}
\noindent\textcolor{red}{Response}:
To avoid the misleading, we modified this sentence to ``For example, assume there is a case with seven master integrals".

%--------------------------------------------
%---------------------------------------------
\begin{center}
\fcolorbox{black}{gray!10}
{
\parbox{1\linewidth}
  {
  Typos and Grammar:
  Page 20, Line 7: only when there exists a $\varphi_J$ share a ...
  }
}
\end{center}
\noindent\textcolor{red}{Response}:
To make it clear, we modified it to ``$\big( \nds \Omega \big)_{IK} = \braket{\dot{\varphi}_I|\varphi_J} \big( \eta^{-1} \big)_{JK}$ could be non-zero only when \textbf{in the summation over $\bm{J}$} there exists a $\varphi_J$ share ($n-1$)-SP or $n$-SP with $\varphi_I$, and $\varphi_J\sim\sim\varphi_K$."


%--------------------------------------------
%---------------------------------------------
\begin{center}
\fcolorbox{black}{gray!10}
{
\parbox{1\linewidth}
  {
  Page 24, last line: Please defne $\nds$ for this case.
  }
}
\end{center}
\noindent\textcolor{red}{Response}:
To make it clear, we added more explanation: ``As a reminder, the $\nds$ here is a formal differential operator with respect to arbitrary one parameter. Since it could act on $c_1$ as well as  $c_2$, we should keep both $\nds \log (z-c_1)$ and $\nds \log (z-c_2)$. Notice that even for $\nds \log z$, we still keep it, because  we regard $\nds \log z$ as $\nds \log (z-c_0)$ with $c_0=0$ and $\nds$ could act on $c_0$ as well."


%--------------------------------------------
%---------------------------------------------
\begin{center}
\fcolorbox{black}{gray!10}
{
\parbox{1\linewidth}
  {
  This example is the same considered in the author’s previous work. Whilst it is instructive to see, how introducing relative twisted cohomology provides the same result as the regulator method, I think the paper could be signifcantly improved by considering an additional new example in relative twisted cohomology, that is not as accessible with the regulator method. This would highlight the signifcance of the paper’s conceptual results.
  }
}
\end{center}
\noindent\textcolor{red}{Response}:
Thanks for this heuristic suggestion. From our current perspective, both relative twisted cohomology and regulator method should work for all cases. The relative twisted cohomology method only take advantages for its simplicity in computation.


%--------------------------------------------
%---------------------------------------------
\begin{center}
\fcolorbox{black}{gray!10}
{
\parbox{1\linewidth}
  {
  Additionally, some comments on how the authors think their ideas could be extended to cases where the appearing forms cannot be written
in $\nd \log$ form, i.e. for Feynman integrals related to elliptic curves or more complicated geometries with exponents in the twist that are not proportional to ε but take e.g. the form $-\frac{1}{2}+\ep$ would be interesting a detailed discussion of this might be beyond the scope of this paper, but one should at least make some comments on this complication and the fact that the case treated in this paper does not cover all Feynman
integrals of interest, which I think is insufciently addressed with the short comment on page 21.
  }
}
\end{center}
\noindent\textcolor{red}{Response}:
Thanks for this valuable suggestion, we added more discussions from eq.(2.97) to the end of this section. In the modification, we try our best to make the $\nd \log$ cases (without elliptic) more clear and roughly outline some basic properties of elliptic cases and the cases beyond elliptic. The reviewers could check the details in the revised version.







%--------------------------------------------
%---------------------------------------------

\noindent\rule[0.25\baselineskip]{\textwidth}{1pt}

% We can also mention that the only difference between our results and former literature is that there is a term has redundant imaginary coefficient $i$ in their results.
We tried our best to improve the manuscript and made some changes in the manuscript. And we appreciate reviewer's warm work earnestly and hope these corrections will clarify points raised by reviewer.

	
% \bibliographystyle{JHEP}
% \bibliography{Reply}
	
	
\end{document}
